% TEMPLATE-MILC-v3.tex - plantilla maestra UNICA de las guias MILC v3.
%
% La usan las cuatro familias de guias, cada una con su driver:
%   · Grados 6-11          -> scripts/build-guias-g11.py
%   · Bebras               -> scripts/build-guias-bebras.py
%   · Semillero            -> scripts/build-guias-semillero.py
%   · Territorio interior  -> scripts/build-guias-territorio-interior.py
%
% Antes habia tres copias casi identicas (~400 lineas iguales, 6 distintas):
% cada mejora habia que aplicarla tres veces, y en la practica no se hacia
% (el motor de recursos vivia solo en dos de ellas). Lo que varia entre
% familias esta parametrizado en 6 placeholders que cada driver llena
% (se nombran aqui SIN su sintaxis real: el builder reemplaza por texto plano,
% tambien dentro de los comentarios, y un valor multilinea rompe el preambulo):
%
%   HEADER_IZQ        encabezado izquierdo de pagina
%   HEADER_DER        encabezado derecho de pagina
%   PORTADA_ETIQUETA  rotulo sobre el numero grande (GRADO/MOMENTO/MODULO)
%   PORTADA_SERIE     linea de serie de la portada
%   PORTADA_ID        identificador de la guia en la portada
%   PORTADA_CIRCULO   el circulito de grado (°), o vacio si no aplica
%   PDF_TITULO        titulo en texto plano (sin \textbf ni comillas TeX) para
%                     los metadatos del PDF (/Title); lo produce lib_xelatex.texto_pdf
%
% Composicion (auditoria 2026-09-03): las bandas de estacion piden espacio
% (\Needspace*) y prohiben el salto justo despues (\nopagebreak), asi no quedan
% huerfanas al pie; las cajas largas (softbox, drawbox, infoband, puentebox) son
% `breakable`, asi una caja de media pagina ya no arrastra todo a la siguiente
% dejando la anterior al 40 %. Todo texto sobre mostaza o verde va en milcNegro
% (blanco sobre #E5B400 da 1,93:1 y sobre #58A923 2,95:1; ninguno pasa WCAG AA).
% El resto de claves las documenta content/guias/_SCHEMA.md. El script valida
% que no queden placeholders sin reemplazar antes de escribir el archivo final.

% Etiquetado PDF (latex-lab): /Lang, estructura y texto alternativo de las
% figuras, para lectores de pantalla. Debe ir ANTES de \documentclass.
\DocumentMetadata{lang=es-CO,tagging=on}
\documentclass[11pt,letterpaper]{article}
% headheight=24pt: la cabecera derecha lleva dos lineas (identidad de la guia
% y estacion actual). headsep se acorta en la misma medida para que el bloque
% de texto quede exactamente donde estaba con headheight=12pt.
\usepackage[margin=2.0cm,headheight=24pt,headsep=13pt]{geometry}
\usepackage[table]{xcolor}
\usepackage{fontspec}
\usepackage[spanish]{babel}
\usepackage{tikz}
% Librerías TikZ para los diagramas de las guías (bloque `recursos:`):
%  · babel  → babel en español vuelve activos caracteres como «>», y TikZ los usa
%             en sus opciones (>=stealth, ->). Sin esta librería, cualquier
%             diagrama con flechas revienta con «\language@active@arg> has an
%             extra }». Debe ir ANTES de que se dibuje nada.
%  · positioning        → sintaxis «below left=2pt and 3pt».
%  · decorations.pathreplacing → llaves ({brace}) para señalar rangos.
\usetikzlibrary{babel,positioning,decorations.pathreplacing}
\usepackage{graphicx}
% Circuitos: el trabajo de electrónica se hace hoy con micro:bit y sus sensores
% integrados (MakeCode, sin cableado), así que no se necesitan esquemáticos.
% Si algún día entran montajes con protoboard, basta descomentar esta línea:
% los diagramas .tex/.tikz declarados en `recursos:` ya se incrustan solos.
% \usepackage{circuitikz}
\usepackage[most]{tcolorbox}
\usepackage{tabularx}
\usepackage{array}
\usepackage{fancyhdr}
\usepackage{titlesec}
\usepackage[document]{ragged2e}  % bandera a la derecha en todo el cuerpo
\usepackage{enumitem}
\usepackage{microtype}
\usepackage{etoolbox}
\usepackage{needspace}
% hyperref al final del preambulo: metadatos del PDF (titulo, autor, idioma,
% familia), marcadores por estacion (\pdfbookmark) y enlaces sin recuadro.
\usepackage[hidelinks,
  pdftitle={Dormir, comer, moverse: la base sobre la que se sostiene el ánimo},
  pdfauthor={PhD. Álvaro Cárdenas Orozco},
  pdfsubject={Territorio interior · Educación socioemocional · Grado 6° · Momento 9 · MILC v3},
  pdflang={es-CO},
  bookmarksopen=true,bookmarksnumbered=false]{hyperref}

% Helvetica Neue no trae la flecha U+2192, asi que cada «→» escrito literalmente
% en el YAML se compilaba a NADA, en silencio (462 flechas en 61 guias: rutas del
% tipo «paleta → tipografia → lockup» salian rotas). Mapeamos el caracter a su
% version matematica, que si tiene glifo. Asi el autor puede seguir escribiendo
% «→» comodamente en el YAML.
\usepackage{newunicodechar}
\newunicodechar{→}{\ensuremath{\rightarrow}}
% Mismo problema con los simbolos de marca y vinetas: el «✓» de «una palomita ✓»
% salia como cuadrito vacio en el PDF de todas las guias que lo usaban.
\usepackage{pifont}
\newunicodechar{✓}{\ding{51}}
\newunicodechar{✗}{\ding{55}}
\newunicodechar{✔}{\ding{52}}
\newunicodechar{•}{\textbullet}
\newunicodechar{≥}{\ensuremath{\geq}}
\newunicodechar{≤}{\ensuremath{\leq}}

\setmainfont{Helvetica Neue}
\setsansfont{Helvetica Neue}
\setlength{\parindent}{0pt}
\setlength{\parskip}{6pt}
% Sin particion de palabras: en 6.º-7.º una palabra cortada por guion al final
% de linea («Comuni-cacion») frena la lectura mucho mas que una linea corta.
\hyphenpenalty=10000
\exhyphenpenalty=10000
\pretolerance=2000
\tolerance=3000
\setlength{\emergencystretch}{3em}
\linespread{1.10}
\renewcommand{\arraystretch}{1.25}
\setlist{leftmargin=5mm,itemsep=1mm,topsep=1mm}
\newcolumntype{Y}{>{\RaggedRight\arraybackslash}X}

\definecolor{milcMagenta}{HTML}{E6007E}
\definecolor{milcVino}{HTML}{5A0038}
\definecolor{milcTurquesa}{HTML}{0093A5}
\definecolor{milcVerde}{HTML}{58A923}
\definecolor{milcMostaza}{HTML}{E5B400}
\definecolor{milcOcre}{HTML}{B86600}
\definecolor{milcNegro}{HTML}{241816}
\definecolor{milcGris}{HTML}{F7F7F4}
\definecolor{milcLinea}{HTML}{E8E2DD}
\definecolor{milcGradePrimary}{HTML}{0D9488}
\definecolor{milcGradeDark}{HTML}{115E59}
\definecolor{milcGradeSoft}{HTML}{CCFBF1}
\definecolor{milcGradeAccent}{HTML}{CFE7FF}
\definecolor{milcGradeText}{HTML}{FFFFFF}
\color{milcNegro}

\pagestyle{fancy}
\fancyhf{}
% Cabecera de dos lineas: identidad de la guia arriba y, a la derecha y en
% gris, la estacion en curso (\markright desde \milcestacion). Antes de la
% primera estacion la segunda linea va vacia; el \strut mantiene la altura
% para que la linea de arriba no baile entre paginas.
\lhead{\small Territorio interior · Educación socioemocional\\[-1pt]{\footnotesize\strut}}
\rhead{\small Grado 6° · Momento 9 · MILC v3\\[-1pt]{\footnotesize\strut\color{milcNegro!65}\rightmark}}
\cfoot{\small\thepage}
\renewcommand{\headrulewidth}{0pt}

% Sobre mostaza (#E5B400) y verde (#58A923) el texto va en milcNegro; sobre el
% resto de la paleta, en blanco. Se usa dentro de fonttitle/fontupper, donde un
% \color posterior manda sobre coltitle.
\newcommand{\milctintasobre}[1]{%
  \ifboolexpr{test{\ifstrequal{#1}{milcMostaza}} or test{\ifstrequal{#1}{milcVerde}}}%
    {\color{milcNegro}}{\color{white}}}

\newtcolorbox{titlebox}[1]{
  enhanced,colback=#1,colframe=#1,arc=7mm,boxrule=0pt,
  left=7mm,right=7mm,top=3mm,bottom=3mm,
  before skip=8pt,after skip=8pt,
  fontupper=\sffamily\bfseries\Large\color{white}
}

\newtcolorbox{softbox}[3]{
  enhanced,breakable,pad at break=3mm,
  colback=#2,colframe=#1,borderline west={4pt}{0pt}{#1},
  arc=2mm,boxrule=.4pt,left=4mm,right=4mm,top=3mm,bottom=3mm,
  before skip=6pt,after skip=8pt,title={#3},coltitle=white,
  colbacktitle=#1,fonttitle=\sffamily\bfseries\milctintasobre{#1},fontupper=\RaggedRight,
  attach boxed title to top left={xshift=3mm,yshift=-2mm},
  boxed title style={arc=2mm, boxrule=0pt}
}

\newtcolorbox{drawbox}[2]{
  enhanced,breakable,pad at break=4mm,
  colback=white,colframe=#1,arc=4mm,boxrule=1pt,
  left=5mm,right=5mm,top=4mm,bottom=4mm,before skip=8pt,after skip=8pt,
  title={#2},coltitle=white,colbacktitle=#1,fonttitle=\sffamily\bfseries\milctintasobre{#1},
  fontupper=\RaggedRight,
  attach boxed title to top left={xshift=4mm,yshift=-2mm},
  boxed title style={arc=3mm, boxrule=0pt}
}

\newtcolorbox{infoband}[1]{
  enhanced,breakable,pad at break=2mm,colback=#1!8,colframe=#1!50,arc=2mm,boxrule=.5pt,
  left=4mm,right=4mm,top=2mm,bottom=2mm,before skip=4pt,after skip=4pt,
  fontupper=\small\RaggedRight
}

% Puente narrativo entre secciones (contrato editorial Conéctate).
% Da continuidad: "Acabas de X. Ahora vas a Y porque Z."
% Visualmente: banda izquierda en color del grado, texto en itálica pequeña.
\newtcolorbox{puentebox}{
  enhanced,breakable,pad at break=2mm,colback=milcGradeSoft,colframe=milcGradeDark,
  borderline west={3pt}{0pt}{milcGradeDark},
  arc=0mm,boxrule=0pt,left=5mm,right=4mm,top=2.5mm,bottom=2.5mm,
  before skip=4pt,after skip=8pt,
  fontupper=\itshape\small\color{milcNegro}\RaggedRight
}
% La flecha va en modo matematico a proposito: Helvetica Neue NO tiene el glifo
% U+2192, asi que \textrightarrow se compilaba a NADA («Missing character: There
% is no → in font Helvetica Neue») y el puente salia sin flecha. En modo
% matematico la flecha viene de la fuente matematica y si se dibuja.
\newcommand{\puente}[1]{%
  \begin{puentebox}%
    \textcolor{milcGradeDark}{$\rightarrow$}\hspace{2mm}#1%
  \end{puentebox}%
}

\newcommand{\linead}[1]{\noindent\color{milcLinea}\rule{#1}{0.5pt}\color{milcNegro}}
\newcommand{\respuesta}[1]{\par\vspace{2mm}\foreach \n in {1,...,#1}{\noindent\color{milcLinea}\rule{\linewidth}{0.5pt}\par\vspace{5mm}}\color{milcNegro}}
\newcommand{\checkbox}{\textcolor{milcGradeDark}{\fbox{\phantom{x}}}\hspace{5pt}}

% ─── Listas aireadas para las guías socioemocionales ────────────────────
% Componer las verificaciones como «\checkbox texto\\» deja el texto que salta
% de línea debajo del cuadrito y apelmaza el bloque. Estos entornos alinean la
% continuación y airean los ítems. Los usa Territorio Interior; cualquier
% familia puede hacerlo.
\newlist{checklista}{itemize}{1}
\setlist[checklista]{
  label=\textcolor{milcGradeDark}{\fbox{\phantom{x}}},
  leftmargin=9mm, labelsep=3.2mm, itemsep=2.4mm,
  topsep=2.5mm, parsep=0pt, partopsep=0pt, before=\RaggedRight
}
\newlist{pasoslista}{enumerate}{1}
\setlist[pasoslista]{
  label=\textcolor{milcGradeDark}{\sffamily\bfseries\arabic*.},
  leftmargin=8mm, labelsep=3mm, itemsep=2.8mm,
  topsep=1mm, parsep=0pt, partopsep=0pt, before=\RaggedRight
}

\newcommand{\phasecard}[4]{
\begin{tcolorbox}[enhanced,colback=#3,colframe=#2,arc=3mm,boxrule=.5pt,height=3.05cm,valign=center,halign=center,left=1.5mm,right=1.5mm,top=2mm,bottom=2mm]
{\sffamily\bfseries\fontsize{22}{24}\selectfont\textcolor{#2}{#1}}\par
{\sffamily\bfseries\small #4}
\end{tcolorbox}
}

% ── Piezas del rediseno 2026 (legibilidad 6.º-11.º) ───────────────────
% Banda de estacion: reemplaza el titlebox macizo. Titulo a la izquierda,
% dato practico (tiempo/modalidad) a la derecha, en una sola linea.
% Nunca huerfana: pide 9 lineas libres (o salta de pagina rellenando con
% \vfill) y prohibe el salto justo despues de la banda. Nueve y no seis:
% la caja partible que sigue necesita ~80pt (titulo adosado, relleno y dos
% lineas) para arrancar en la misma pagina; con menos, tcolorbox la manda
% entera a la siguiente y la banda se queda sola. El penalty 10000 va
% en `after`, ANTES del espacio posterior: un `after skip` (aunque sea 0pt)
% es un \vskip tras la caja, y ese glue es un punto de corte legal que un
% \nopagebreak escrito despues de \end{tcolorbox} ya no cierra. Cada banda
% deja marcador PDF y actualiza la cabecera derecha.
\newcounter{milcestacion}
\newcommand{\milcestacion}[2]{%
\Needspace*{9\baselineskip}%
\stepcounter{milcestacion}%
\pdfbookmark[1]{#2}{milcestacion\themilcestacion}%
\markright{#2}%
\begin{tcolorbox}[enhanced,colback=#1,colframe=#1,arc=2mm,boxrule=0pt,
  left=5mm,right=5mm,top=2.6mm,bottom=2.6mm,before skip=11pt,
  after={\par\nopagebreak\vskip7pt}]
{\sffamily\bfseries\fontsize{15}{18}\selectfont\milctintasobre{#1}#2}
\end{tcolorbox}}

% Tarjeta de paso numerada: sustituye las tablas «Pilar / Aplicacion», que
% eran formato de consulta docente, no de estudiante. El numero va en un
% circulo sobre el borde para que la secuencia se lea de un vistazo.
\newcommand{\milcpaso}[3]{%
\Needspace{4\baselineskip}%
\begin{tcolorbox}[enhanced,colback=white,colframe=milcLinea,arc=1.5mm,boxrule=.9pt,
  left=14mm,right=4mm,top=3mm,bottom=3mm,before skip=4pt,after skip=4pt,
  fontupper=\RaggedRight,
  overlay={\node[anchor=north west,circle,fill=#2,text=white,inner sep=0pt,
    minimum size=7.4mm,font=\sffamily\bfseries\small]
    at ([xshift=3.6mm,yshift=-3.2mm]frame.north west) {#1};}]
#3
\end{tcolorbox}}

% Casilla marcable de tamano util con lapiz (la anterior era \fbox{\phantom{x}},
% de alto variable segun el texto que la seguia).
\newcommand{\milcmarca}{\framebox[3.8mm]{\rule{0pt}{3.4mm}}}

% Fila marcable de lista suelta (checks de la estacion 1).
\newcommand{\milccheckrow}[1]{%
\par\vspace{1.8mm}\noindent
\makebox[6.4mm][l]{\raisebox{-.55mm}{\textcolor{milcNegro}{\milcmarca}}}%
\parbox[t]{\dimexpr\linewidth-6.4mm\relax}{\RaggedRight #1}\par}

% Celda marcable dentro de la rubrica (tabla de autoevaluacion). Misma
% sangria francesa, pero pensada para vivir dentro de una columna X.
\newcommand{\milcopcion}[1]{%
\makebox[5.2mm][l]{\raisebox{-.55mm}{\textcolor{milcNegro}{\milcmarca}}}%
\parbox[t]{\dimexpr\linewidth-5.2mm\relax}{\RaggedRight #1}}

% ── Recursos visuales de la guía ──────────────────────────────────────
% Declarados en el bloque `recursos:` del YAML y emitidos por el builder.
% \guiaFigura   → imagen rasterizada o PDF (foto, carta celeste, montaje).
% \guiaDiagrama → fuente TikZ/circuitikz (.tex) incrustada como vector:
%                 es la vía para circuitos y esquemáticos editables.
% [#1] = texto alternativo (alt) · #2 = ruta relativa al .tex · #3 = pie
% El alt llega al PDF etiquetado (/Alt) y lo lee el lector de pantalla.
\newcommand{\guiaFigura}[3][]{%
\begin{center}
\includegraphics[width=0.88\linewidth,height=0.38\textheight,keepaspectratio,alt={#1}]{#2}\\[1.5mm]
{\footnotesize\itshape\textcolor{milcNegro!75}{#3}}
\end{center}
}

\newcommand{\guiaDiagrama}[2]{%
\begin{center}
\input{#1}\\[1.5mm]
{\footnotesize\itshape\textcolor{milcNegro!75}{#2}}
\end{center}
}

\begin{document}
\thispagestyle{empty}

% ── Portada ───────────────────────────────────────────────────────────
% Antes: un rectangulo de color a pagina completa. Se veia bien en pantalla,
% pero en la fotocopiadora del colegio se comia el toner y salia gris sucio.
% Ahora la banda ocupa el tercio superior y el resto va en blanco.
\begin{tikzpicture}[remember picture,overlay]
  \path[fill=milcGradePrimary]
    (current page.north west) rectangle ([yshift=-10.6cm]current page.north east);

  \node[anchor=north west,text=milcGradeText] at ([xshift=2.0cm,yshift=-1.55cm]current page.north west)
    {\sffamily\bfseries\fontsize{12}{14}\selectfont MOMENTO};
  \node[anchor=north west,text=milcGradeText,opacity=.85] at ([xshift=2.0cm,yshift=-2.35cm]current page.north west)
    {\sffamily\bfseries\fontsize{8.5}{10}\selectfont TERRITORIO INTERIOR · SOCIOEMOCIONAL};
  \node[anchor=north east,text=milcGradeText,opacity=.85,align=right] at ([xshift=-2.0cm,yshift=-1.55cm]current page.north east)
    {\sffamily\bfseries\fontsize{9}{12}\selectfont I.E. Sor María Juliana\\Cartago, Valle del Cauca};

  \node[anchor=west,text=milcGradeText] (gradonum) at ([xshift=1.85cm,yshift=-7.1cm]current page.north west)
    {\sffamily\bfseries\fontsize{112}{112}\selectfont 9};
  % El circulito de grado (°) solo aplica a las guias de grado («11°»); en
  % Bebras y Semillero el numero grande es un momento/modulo, no un grado.
  % Se posiciona relativo al nodo `gradonum`, no en coordenadas absolutas.
  

  \node[anchor=west,text=milcGradeText,text width=11.4cm,align=left] at ([xshift=1.5cm,yshift=0.5cm]gradonum.east)
    {
     {\sffamily\bfseries\fontsize{9}{11}\selectfont\MakeUppercase{Territorio interior · Socioemocional}}\\[3mm]
     {\sffamily\bfseries\fontsize{21}{25}\selectfont\setlength{\baselineskip}{27pt}Dormir, comer, moverse\\[4pt]la base del ánimo\\[4pt]no es un lujo}\\[3.5mm]
     {\sffamily\bfseries\fontsize{15}{18}\selectfont Grado 6° · Momento 9}
    };
\end{tikzpicture}

\vspace*{9.4cm}
\vfill

{\sffamily\bfseries\fontsize{8}{10}\selectfont\textcolor{milcGradeDark}{LO QUE TE LLEVAS HOY}}

\begin{tcolorbox}[enhanced,colback=white,colframe=milcNegro,arc=2mm,boxrule=1.6pt,
  left=5mm,right=5mm,top=4mm,bottom=4mm,before skip=3pt,after skip=12pt,
  fontupper=\RaggedRight\fontsize{11}{15.5}\selectfont]
Tu calendario propio: el registro de una semana de tus tres bases —sueño, comidas y movimiento— cruzado con tu ánimo del día, y un solo cambio elegido, del tamaño de una semana.
\end{tcolorbox}

\begin{tcolorbox}[enhanced,colback=milcGris,colframe=milcGris,arc=2mm,boxrule=0pt,
  left=5mm,right=5mm,top=3.5mm,bottom=3.5mm,after skip=12pt]
{\sffamily\bfseries\fontsize{8}{10}\selectfont\textcolor{milcNegro!55}{ESTA GUÍA ES DE}}\\[3mm]
\begin{tabularx}{\linewidth}{X p{3.2cm} p{3.2cm}}
\linead{\linewidth} & \linead{3.0cm} & \linead{3.0cm}\\[-1mm]
{\fontsize{8.5}{10}\selectfont\textcolor{milcNegro!55}{Nombre completo}} &
{\fontsize{8.5}{10}\selectfont\textcolor{milcNegro!55}{Grupo}} &
{\fontsize{8.5}{10}\selectfont\textcolor{milcNegro!55}{Fecha}}\\
\end{tabularx}
\end{tcolorbox}

\vfill

\noindent\textcolor{milcLinea}{\rule{\linewidth}{1pt}}\\[2mm]
{\sffamily\bfseries\fontsize{9.5}{12}\selectfont Autor: Dr. Álvaro Cárdenas Orozco}\\[1mm]
{\sffamily\fontsize{8.5}{11}\selectfont\textcolor{milcNegro!55}{Metodología MILC · Apertura ancestral · Cuidado del cuerpo · Triángulo Dussel-Estoicismo-Floridi}}

\newpage

\section*{Apertura: Dormir, comer, moverse: la base sobre la que se sostiene el ánimo}

\begin{softbox}{milcGradeDark}{milcGradeSoft}{Saber ancestral}
Colombia es uno de los pocos países del mundo con \textbf{dos cosechas de café al año}, y eso no es mérito de nadie: es la posición sobre la línea del ecuador y el régimen de lluvias. Está la \textbf{cosecha principal} ---entre septiembre y diciembre en el centro y el norte--- y una segunda más pequeña, la \textbf{mitaca}, también llamada \textbf{traviesa}, entre abril y junio. El Valle del Cauca está entre los departamentos que tienen traviesa. \textbf{Y aquí está lo que enseña el calendario:} la finca \emph{no produce parejo todo el año}. Hay semanas de recolección en que se trabaja de sol a sol y toda la vereda se vuelca al cafetal; y hay meses en que no hay casi nada que recoger, porque la mata está floreciendo y llenando el grano. \textbf{Ese tiempo no es tiempo perdido: es donde se prepara la cosecha siguiente.} La cosecha depende de la floración, de la lluvia y de la luz ---cosas que ningún caficultor puede apurar---. \emph{El que le exige a un cafetal producir al máximo todo el año no obtiene más café: obtiene una mata agotada.} Tu cuerpo funciona igual, y a los once años vale la pena saberlo antes de creer lo contrario.
\end{softbox}

\begin{softbox}{milcTurquesa}{milcGris}{Saber contemporáneo}
Casi todo lo que aprendiste en los momentos anteriores ---nombrar, frenar, respirar--- \textbf{funciona peor con el cuerpo en el piso}. No es una cuestión de carácter: cuando hay sueño de menos, el mismo comentario duele más, la misma tarea parece imposible y la pausa no alcanza. Por eso este momento no es un consejo de salud pegado al final del programa: \textbf{es la base sobre la que se sostiene todo lo demás}. Los datos son claros. La \textbf{Academia Americana de Medicina del Sueño} revisó la evidencia y recomienda, para \textbf{niños de 6 a 12 años, entre 9 y 12 horas} de sueño por cada 24, y para \textbf{adolescentes de 13 a 18 años, entre 8 y 10} (Paruthi y otros, 2016). \emph{Compáralo con lo que dormiste anoche, sin dramatizar y sin mentirte.} Súmale las otras dos bases: \textbf{comer} algo de verdad en el desayuno y en el descanso ---el cerebro trabaja con lo que le llegue--- y \textbf{moverse}, aunque sea caminar. \textbf{No son tres consejos sueltos: son las tres patas de la misma silla}, y cuando falta una, se siente en el ánimo antes que en el cuerpo.
\end{softbox}

\begin{softbox}{milcMostaza}{milcGris}{Pregunta-puente}
\textit{El caficultor sabe que la mata no da si no descansa, y no lo considera pereza: lo considera parte del oficio. \textbf{¿Y si tu cansancio de estas semanas no fuera falta de ganas}\ldots sino una cosecha que llevas exigiéndote sin haber tenido traviesa?}
\end{softbox}

\begin{softbox}{milcVerde}{milcGris}{Saber y saber hacer}
Al terminar podrás: (1) \textbf{identificar} cómo estuvieron tus tres bases ---sueño, comidas y movimiento--- en los últimos días, con datos; (2) \textbf{explicar} por qué el cuerpo funciona por ciclos y cuánto sueño necesita alguien de tu edad, según la evidencia; (3) \textbf{analizar} tu propia semana buscando la relación entre esas bases y tu ánimo; (4) \textbf{decidir} un solo cambio, del tamaño de una semana, y saber cómo comprobarás si sirvió.
\end{softbox}



\small
\begin{tabularx}{\textwidth}{>{\bfseries}p{3.0cm}Y}
\rowcolor{milcGradePrimary!12}Autor & Dr. Álvaro Cárdenas Orozco\\
DBA & Comprende los fenómenos que afectan a su territorio, regula sus emociones ante ellos y acompaña a otros con empatía (transversal 6°--11°, articulado con la política «Escuela, territorio de vida» del MEN, Resolución 006519 de 2025, y la Ley 1620 de 2013)\\
Referentes & Primera ayuda psicológica (OMS, 2011/2012) · Directrices IASC (2007) · USGS y Servicio Geológico Colombiano · Pueblo nasa y la minga (CRIC) · Dussel · Estoicismo · Floridi\\
Producto final & Tu calendario propio: el registro de una semana de tus tres bases —sueño, comidas y movimiento— cruzado con tu ánimo del día, y un solo cambio elegido, del tamaño de una semana.\\
\end{tabularx}
\normalsize

\puente{Llevas cuatro momentos aprendiendo a manejar lo que sientes. Hoy vas a lo que está debajo de todo eso: \textbf{si el cuerpo no está, las herramientas no alcanzan}.}

\milcestacion{milcGradeDark}{Ruta MILC: así vamos a aprender}
\begin{tabularx}{\textwidth}{>{\centering\arraybackslash}X>{\centering\arraybackslash}X>{\centering\arraybackslash}X>{\centering\arraybackslash}X}
\phasecard{1}{milcVerde}{milcGris}{Escucha\\[-1mm]\small Observo, escucho y pregunto.} &
\phasecard{2}{milcTurquesa}{milcGris}{Sistematización\\[-1mm]\small Observo y relaciono.} &
\phasecard{3}{milcMagenta}{milcGris}{Praxis\\[-1mm]\small Construyo y aplico.} &
\phasecard{4}{milcVino}{milcGris}{Evaluación\\[-1mm]\small Reflexión liberadora.}
\end{tabularx}


\puente{Primero los datos de tus últimos días. Sin adornarlos y sin juzgarte: es una medición, no una acusación.}

\milcestacion{milcVerde}{Estación 1 de 3 · Escucha}

\begin{softbox}{milcVerde}{milcGris}{Escena para escuchar}
\textbf{Actividad 1 · IDENTIFICA --- Mis tres bases, ayer y anteayer.} \textbf{Nadie va a leer esta página.} Para \textbf{cada uno} de los dos últimos días, anota cinco datos: \textbf{a qué hora te dormiste} y \textbf{a qué hora te levantaste} (y saca las horas); \textbf{si desayunaste} y qué; \textbf{si comiste algo en el descanso}; \textbf{cuánto te moviste} ---no ejercicio: moverte, caminar, jugar, subir escaleras---; y \textbf{cómo estuvo tu ánimo}, del 1 al 10, con una palabra de tu muestrario. \textbf{Después, una pregunta:} de los dos días, ¿cuál estuvo mejor de ánimo? Míralo al lado de las otras cuatro columnas. \emph{No saques conclusiones todavía con dos días ---eso es muy poco para concluir nada---, pero fíjate si algo salta a la vista. Casi siempre salta.}
\end{softbox}

{\sffamily\bfseries\fontsize{8}{10}\selectfont\textcolor{milcNegro!55}{MARCA CUANDO LO HAYAS HECHO}}
\milccheckrow{Anotaste, para los dos días, la hora de acostarte y de levantarte, con el total de horas.}
\milccheckrow{Registraste desayuno, algo en el descanso y cuánto te moviste, sin adornarlo.}
\milccheckrow{Le pusiste número y palabra a tu ánimo de cada día, y los comparaste.}

\begin{infoband}{milcVerde}


\textbf{En tu cuaderno (Actividad 1 · IDENTIFICA).} Título: «Actividad 1. Mis tres bases». \textbf{Formato:} una tabla de dos filas (los dos días) y cinco columnas (horas de sueño · desayuno · descanso · movimiento · ánimo 1-10 y palabra). \textbf{Extensión:} la tabla y dos renglones de observación. \textbf{Sabes que terminaste cuando} los datos son \textbf{los reales}, no los que te gustaría tener. Esta página es tuya: \textbf{nadie más tiene que leerla si tú no quieres}.
\end{infoband}

\puente{Ya tienes tus números. Ahora entiende por qué el cuerpo va en ciclos y cuánto necesita de verdad.}

\milcestacion{milcTurquesa}{Fase 2 · Sistematización: entender el ciclo}

\begin{softbox}{milcTurquesa}{milcGris}{Cuatro claves sobre la base física del ánimo}
Cuatro claves para entender por qué esto no es un sermón de salud, sino la base de todo lo demás que has trabajado este año.
\end{softbox}

\milcpaso{1}{milcTurquesa}{\textbf{El cuerpo va en ciclos, como el cafetal.} No estás hecho para rendir igual todos los días ni a todas las horas: hay franjas del día en que tu cabeza trabaja mejor y otras en que no, y semanas de más carga y semanas de menos. \textbf{Reconocer tu ciclo no es excusarte: es planear con lo que hay.} Si sabes que a las nueve de la noche ya no te entra nada, la solución no es tomar café: es mover esa tarea a la mañana. \emph{El caficultor no le grita a la mata para que florezca antes; lee el ciclo y se organiza alrededor de él.}}
\milcpaso{2}{milcTurquesa}{\textbf{Cuánto sueño hace falta, con número.} La Academia Americana de Medicina del Sueño lo revisó y lo dejó claro: de \textbf{6 a 12 años, entre 9 y 12 horas} por cada 24; de \textbf{13 a 18, entre 8 y 10} (Paruthi y otros, 2016). \textbf{No son horas «ideales»: son las que el cuerpo pide para funcionar.} Y ojo con la trampa más común a tu edad: dormir poco entre semana y \emph{recuperar} el fin de semana funciona a medias ---ayuda algo, pero deja el resto de la semana en deuda---. \emph{Si de tus últimos dos días ninguno llegó a nueve horas, ahí tienes una explicación de más cosas de las que crees.}}
\milcpaso{3}{milcTurquesa}{\textbf{Comer y moverse no son adornos.} El cerebro trabaja con lo que le llegue: pasar la mañana sin nada hace que a la tercera hora todo cueste más ---incluido aguantarse las ganas de responder mal---. Y el movimiento no es «hacer ejercicio»: es \textbf{no pasar el día sentado}. Caminar al colegio, jugar en el descanso, subir escaleras, bailar en la casa. \emph{Lo que casi nadie te dice es que el efecto se siente primero en el ánimo y solo después en el cuerpo:} el mal genio de las once suele ser hambre, y el bajón de la tarde suele ser quietud.}
\milcpaso{4}{milcTurquesa}{\textbf{La cosecha no es todo el año, y la traviesa no es pereza.} Habrá semanas de cosecha ---exámenes, entregas, un partido importante--- en que se duerme menos y se aprieta. \textbf{Eso está bien si es una temporada y no la vida entera.} Lo que revienta a la gente no es la semana dura: es no tener nunca traviesa. \emph{Programa tu descanso como programas tu trabajo:} después de la semana de exámenes, un día de dormir sin alarma; después de la entrega, una tarde sin nada. Y si te toca una cosecha larga, avísale a alguien: eso es lo del momento siguiente.}

\begin{drawbox}{milcTurquesa}{Las tres bases y qué se ve cuando falla cada una}
\begin{pasoslista}
\item \textbf{Sueño (9 a 12 horas a tu edad).} Cuando falta: todo irrita más, la cabeza se dispersa, la pausa no alcanza y las emociones suben más rápido.
\item \textbf{Comida (desayuno y algo en el descanso).} Cuando falta: bajón a media mañana, mal genio que uno cree que es de carácter, dificultad para sostener la atención.
\item \textbf{Movimiento (moverse, no entrenar).} Cuando falta: ánimo plano, sueño peor en la noche, cuerpo tenso sin causa clara.
\item \textbf{Y la regla del ciclo:} después de una temporada de cosecha, va una de traviesa. Si nunca llega, el problema no es tu ánimo: es tu calendario.
\end{pasoslista}
\end{drawbox}

\begin{softbox}{milcMostaza}{milcGris}{Errores comunes y su corrección}
\textbf{«Yo con cinco horas funciono»:} funcionar no es lo mismo que estar bien; el que duerme poco es mal juez de cuánto le está costando. \textbf{Recuperar todo el fin de semana:} ayuda algo, pero no cierra la deuda de la semana. \textbf{Confundir descanso con pantalla:} media hora de videos no descansa igual que dormir o caminar ---eso lo viste en el momento anterior---. \textbf{Creer que el ánimo es solo carácter:} el mal genio de las once suele ser hambre. \textbf{Cambiar cinco cosas a la vez:} no se sostienen y no sabrás cuál sirvió. \textbf{Convertir esto en culpa:} en muchas casas la hora de dormir o el desayuno no dependen del todo de uno; se trabaja con lo que sí depende.
\end{softbox}

\begin{infoband}{milcTurquesa}


\textbf{En tu cuaderno (Actividad 2 · EXPLICA).} Título: «Actividad 2. Por qué el cuerpo manda». \textbf{Formato:} las tres bases con \textbf{qué te pasa a ti} cuando falla cada una (una línea por base, con un ejemplo tuyo real) + cuántas horas dormiste anoche frente a las 9-12 recomendadas. \textbf{Extensión:} media página. \textbf{Sabes que terminaste cuando} puedes nombrar \textbf{una discusión o un mal rato} de este mes que probablemente empezó por sueño o por hambre.
\end{infoband}

\puente{Ya lo entiendes. Es hora de llevar el registro de una semana, que es la única manera de ver la relación con tus propios ojos.}

\milcestacion{milcMagenta}{Fase 3 · Praxis: armar tu calendario}

\begin{softbox}{milcMagenta}{milcGris}{Cómo se lleva el registro de una semana}
Dos días no alcanzan para ver un ciclo: una semana sí. Hoy armas el registro y eliges \textbf{un} cambio ---uno solo--- del tamaño de siete días.
\end{softbox}

\milcpaso{1}{milcMagenta}{\textbf{Las tres columnas (dibuja la tabla hoy).} Siete filas ---una por día--- y estas columnas: \emph{horas de sueño · desayuno sí/no · algo en el descanso sí/no · movimiento (minutos aproximados) · ánimo del 1 al 10 · palabra del muestrario}. \textbf{Dibújala hoy en clase}, no en la casa: una tabla que no existe no se llena.}
\milcpaso{2}{milcMagenta}{\textbf{El ánimo del día, con número y palabra.} El número sirve para comparar; la palabra ---la del momento 2--- sirve para entender. \emph{«6, aburrido» y «6, irritado» son dos días muy distintos, y esa diferencia es la que después explica todo.} Anótalo siempre a la misma hora, de noche, para que sea comparable.}
\milcpaso{3}{milcMagenta}{\textbf{Buscar el cruce (la semana entrante, en clase).} Con la tabla llena se hace lo interesante: \textbf{ordenar los días de mejor a peor ánimo} y mirar qué tenían en común los tres mejores y los tres peores. \emph{No hace falta estadística: si tus tres peores días durmieron menos de siete horas, ya tienes tu respuesta ---y la sacaste tú, de tus datos, que es distinto a que te la digan---.}}
\milcpaso{4}{milcMagenta}{\textbf{Un solo cambio (elígelo hoy).} Del tamaño de una semana, concreto y \textbf{que dependa de ti}: \emph{«apago el celular a las 9:30 y lo dejo cargando afuera»} ---que además engancha con tus reglas del momento anterior---; \emph{«desayuno aunque sea un huevo o un pan con queso»}; \emph{«me devuelvo caminando del colegio tres días»}. \textbf{Uno.} Y escribe cómo sabrás si sirvió: \emph{«si el promedio de mi ánimo de la semana sube de \_\_\_\_ a \_\_\_\_»}.}

\begin{drawbox}{milcMagenta}{Antes de cerrar tu calendario}
\begin{checklista}
\item La tabla de siete filas está \textbf{dibujada}, con sus seis columnas, hoy y en clase.
\item Sabes a qué hora vas a anotar el ánimo cada día, para que los días sean comparables.
\item Tu cambio es \textbf{uno solo}, concreto y depende de ti (no de que otro haga algo).
\item Escribiste cómo vas a saber si sirvió, con un número, no con «me sentí mejor».
\item Anotaste cuántas horas dormiste anoche y cuántas recomienda la evidencia para tu edad.
\item Si algo de tus bases no depende de ti ---la hora de la casa, si hay desayuno---, lo escribiste aparte para hablarlo, no para cargarlo.
\end{checklista}
\end{drawbox}

\begin{softbox}{milcMagenta}{milcGris}{Plantilla de trabajo}
\textbf{Plantilla del calendario (siete filas).} \emph{Día} · \emph{Horas de sueño} · \emph{Desayuno} · \emph{Descanso} · \emph{Movimiento (min)} · \emph{Ánimo 1-10} · \emph{Palabra}. \\[1mm]
\textbf{Mi cambio de la semana.} \emph{«Voy a \_\_\_\_, todos los días, y sabré que sirvió si \_\_\_\_ (número)».} \\[1mm]
\textbf{Lo que no depende de mí:} \_\_\_\_ \emph{(para hablarlo con \_\_\_\_, no para cargarlo solo).}
\end{softbox}

\begin{infoband}{milcMagenta}


\textbf{En tu cuaderno (Actividad 3 · APLICA).} Título: «Actividad 3. Mi calendario de la semana». \textbf{Formato:} la tabla de siete filas (vacía hoy) + el cambio elegido con su forma de verificarlo + la lista de lo que no depende de ti. \textbf{Extensión:} una página. \textbf{Sabes que terminaste cuando} tu cambio es tan pequeño que \textbf{te da un poco de risa} ---esos son los que se sostienen---.
\end{infoband}

\puente{El producto es ese registro y \textbf{un} cambio. \emph{Cambiar cinco cosas a la vez es la forma más segura de no sostener ninguna ---y de no saber cuál sirvió---.}}

\milcestacion{milcMostaza}{Producto de la sesión}
\textbf{Calendario de una semana: tres bases, ánimo diario y un solo cambio verificable}.
\begin{softbox}{milcMostaza}{milcGris}{Criterios de calidad}
(1) La tabla tiene siete filas y las seis columnas, y quedó dibujada en clase. (2) El ánimo se registra con número y con palabra del muestrario, siempre a la misma hora. (3) El cambio elegido es uno solo, concreto y depende de quien lo escribe. (4) Está escrito cómo se comprobará si sirvió, con una medida. (5) Lo que no depende del estudiante está anotado aparte, para conversarlo con un adulto.
\end{softbox}

\puente{Antes de cerrar, mira si tu cambio sobrevive a una semana con examen. Si solo funciona en una semana tranquila, todavía no es tuyo.}

\milcestacion{milcVino}{Evaluación liberadora · cinco dimensiones}

\begin{softbox}{milcVino}{milcGris}{Sentido de la evaluación}
Evaluar no es solo revisar una nota. Es reconocer qué aprendí, qué cambió en mí, qué emociones manejé, cómo me relaciono con la ciudad y la comunidad, y qué le dejaré a quien viene detrás.
\end{softbox}

\textbf{Mi reflexión liberadora en cinco dimensiones.} Completa cada frase iniciada:

\small
\begin{tabularx}{\textwidth}{>{\bfseries\RaggedRight\arraybackslash}p{3.4cm}Y>{\RaggedRight\arraybackslash}p{4.6cm}}
\rowcolor{milcVino}\color{white}Dimensión & \color{white}\bfseries Mi respuesta (frame starter) & \color{white}\bfseries Algo que descubrí\\
1. Personal & Lo que más cambió en mí fue \linead{6.5cm} & \linead{4.2cm}\\
2. Emocional & La emoción más fuerte fue \linead{2.5cm} y la regulé \linead{3cm} & \linead{4.2cm}\\
3. Ciudadana & En democracia esto importa porque \linead{6cm} & \linead{4.2cm}\\
4. Local (Cartago) & En mi barrio sería útil para \linead{6.5cm} & \linead{4.2cm}\\
5. Intergeneracional & A mi abuela/abuelo le diría \linead{6.5cm} & \linead{4.2cm}\\
\end{tabularx}
\normalsize

\begin{infoband}{milcVino}
\textbf{Para la próxima sustentación.} Vuelve a esta tabla en seis meses. Verás que las respuestas cambian — no porque lo aprendido cambie, sino porque tú cambias. La evaluación liberadora es una conversación contigo a través del tiempo.
\end{infoband}


\puente{Cerramos con tres voces: el cuerpo del que no puede elegir, lo que sí depende de ti, y qué significa cuidar el organismo que también somos.}

\milcestacion{milcGradeDark}{Triángulo de pensamiento · cierre filosófico}

\begin{softbox}{milcVino}{milcGris}{Dussel · lente del nosotros}
\textit{«El otro se revela realmente como otro\ldots como el pobre, el oprimido; el que, a la vera del camino, fuera del sistema, muestra su rostro sufriente y sin embargo desafiante: «¡Tengo hambre!, ¡tengo derecho a comer!»»}

{\footnotesize\itshape\color{milcNegro!70}--- Enrique Dussel · Filosofía de la liberación (1977)}

Dussel no eligió esa frase por casualidad: empieza por el \textbf{hambre}, que es lo más corporal que hay. Y eso obliga a decir algo con cuidado en un salón de sexto: \textbf{no todos llegan igual}. Hay quien no desayunó, quien duerme con tres personas en el mismo cuarto, quien trabajó anoche. \emph{Cuando alguien está irritable a las once, la pregunta honesta no es «¿por qué eres así?» sino «¿comiste?».} \textbf{Y al revés:} si lo que te falta no depende de ti, eso no se aguanta en silencio ---se le dice a un adulto, que es de lo que trata el momento siguiente---.

\textbf{¿A quién de mi salón he juzgado por su mal genio sin preguntarme si durmió o si comió?}
\end{softbox}
\linead{\linewidth}\\[1mm]
\linead{\linewidth}

\begin{softbox}{milcMostaza}{milcGris}{Estoicismo · Epicteto · Enquiridión, 1 (c. 125 d.C.; trad. desde E. Carter) · cuidado interior}
\textit{«Algunas cosas están en nuestro poder y otras no. Están en nuestro poder la opinión, el impulso, el deseo, el rechazo; en una palabra, todo aquello que es obra nuestra.»}

{\footnotesize\itshape\color{milcNegro!70}--- Epicteto · Enquiridión, 1 (c. 125 d.C.; trad. desde E. Carter)}

A los once años \textbf{muchas de estas bases no dependen del todo de ti}: la hora en que se apaga la casa, si hay desayuno, cuánto ruido hay para dormir. Fingir que sí depende todo de ti sería mentirte y cargarte con una culpa que no es tuya. \textbf{Pero algo siempre queda en tu poder}, y ahí es donde se trabaja: la hora en que sueltas el teléfono, si te devuelves caminando, si te tomas el agua, si avisas cuando llevas tres noches sin dormir. \emph{Separar las dos columnas ---como en el momento 3--- es lo que convierte esto en un plan y no en un reclamo.}

\textbf{De mis tres bases, ¿cuál parte depende de mí y cuál tendría que hablar con alguien de mi casa?}
\end{softbox}
\linead{\linewidth}\\[1mm]
\linead{\linewidth}

\begin{softbox}{milcTurquesa}{milcGris}{Floridi · lente de la infoesfera}
\textit{«Somos organismos informacionales ---inforgs--- que compartimos un entorno global hecho, en última instancia, de información: la infoesfera.»}

{\footnotesize\itshape\color{milcNegro!70}--- Luciano Floridi · La cuarta revolución (2014)}

Floridi dice que somos \emph{organismos} informacionales, y conviene subrayar la primera palabra: \textbf{organismos}. Se puede vivir conectado a un entorno hecho de información y seguir siendo un cuerpo que necesita nueve horas de sueño, comida y movimiento. \emph{El entorno digital no tiene horario ni cansancio ---la lista no se acaba, la conversación sigue a las dos de la mañana---, y por eso hay que ponerle nosotros el límite que él no trae.} \textbf{Tu cuerpo no se actualiza a la velocidad de tu teléfono.}

\textbf{¿Cuántas de mis horas de sueño se las está llevando algo que podría esperar a mañana?}
\end{softbox}
\linead{\linewidth}\\[1mm]
\linead{\linewidth}

\begin{infoband}{milcNegro}
\textbf{Nota para el docente.} Las tres son citas textuales y llevan su referencia debajo. Puedes remitir a la obra en clase.
\end{infoband}

\milcestacion{milcVino}{Mi autoevaluación · marca una casilla por fila}

{\small
\setlength{\extrarowheight}{2.2pt}
\begin{tabularx}{\textwidth}{>{\RaggedRight\arraybackslash\bfseries}p{3.0cm}YYY}
\rowcolor{milcVino}\color{white}\bfseries Criterio & \color{white}\bfseries Lo logré & \color{white}\bfseries En proceso & \color{white}\bfseries Necesito apoyo\\[.6mm]
Comprensión del tema & \milcopcion{Explico las ideas con mis palabras.} & \milcopcion{Reconozco las ideas, falta precisión.} & \milcopcion{Necesito repasar lo básico.}\\[.8mm]
\rowcolor{milcGris}Lectura del propio ciclo & \milcopcion{Relaciono mi ánimo con lo que dormí, comí y me moví, con datos de mi semana.} & \milcopcion{Registro los datos pero todavía no encuentro la relación.} & \milcopcion{Creo que mi ánimo no tiene nada que ver con mi cuerpo.}\\[.8mm]
Aplicación y producto & \milcopcion{Mi producto cumple los criterios.} & \milcopcion{Está armado, con ajustes pendientes.} & \milcopcion{Aún está en construcción.}\\[.8mm]
\rowcolor{milcGris}Reflexión 5 dimensiones & \milcopcion{Respondí cada una con honestidad.} & \milcopcion{Respondí algunas con detalle.} & \milcopcion{Mis respuestas son muy generales.}\\[.8mm]
Triángulo de pensamiento & \milcopcion{Conecté las 3 voces con mi trabajo.} & \milcopcion{Conecté al menos una.} & \milcopcion{No las relaciono con mi proceso.}\\[.8mm]
\end{tabularx}}

\vspace{3mm}
\textbf{Hoy entendí que:}\\[1mm]
\linead{\linewidth}\\[1mm]
\linead{\linewidth}

\textbf{Mi compromiso para la próxima sesión es:}\\[1mm]
\linead{\linewidth}\\[1mm]
\linead{\linewidth}

\begin{softbox}{milcMostaza}{milcGris}{Cita de despedida · Marco Aurelio}
\textit{``El obstáculo en el camino se vuelve el camino. Lo que estorba al camino se vuelve camino.''}\\[1mm]
Cada error de hoy, bien observado, se convierte en pieza del camino que pisarás mañana.
\end{softbox}

\small
\begin{tabularx}{\textwidth}{>{\bfseries}p{3.6cm}Y}
\rowcolor{milcGradeDark!12}Nombre completo: & \linead{10cm}\\
Fecha: & \linead{4cm} \quad Curso/grupo: \linead{4cm}\\
Mi firma: & \linead{9cm}\\
\end{tabularx}
\normalsize

\begin{infoband}{milcGradeDark}
\textbf{Para la próxima sesión.} Dos cosas. \textbf{Primero, llena las siete filas} y trae la tabla completa, con tu cambio aplicado y la comprobación hecha ---subió, bajó o quedó igual: los tres resultados sirven---. \textbf{Segundo, pregunta por el calendario de tu familia:} averigua con alguien mayor \textbf{cuándo era la época de más trabajo} en su vida ---la cosecha, el fin de mes, la temporada de ventas, la siembra--- y \textbf{qué hacían cuando pasaba}. Si vienen de finca, pregunta directamente por la cosecha y la traviesa: cómo se organizaba la casa en cada una. \textbf{Trae:} tu tabla llena y lo que te contaron. \emph{Vas a descubrir que la idea de alternar temporadas de mucho y de poco no la inventó nadie en un libro de autoayuda: es la manera en que ha trabajado la gente que vive del campo desde siempre.}
\end{infoband}

\end{document}
